%% kapa: exam=1 round=20 year=2009 period=2 form=A subject=accounting subject_ko=회계학 seq=26 number=26
%% source: Q-Net 감정평가사 기출문제 (cst00309) artlSeq=5024178 file=extracted/1008/5024178/2009년 제20회 감정평가사 1차시험 기출문제/2교시.hwp
%% conversion: kordoc 4.24.51 -> tools/gen1.py (auto); review=proofread-sonnet; flags=none
\begin{question}{26}
(주)한국은 20×7년 초에 본사 사옥을 신축하기 위해서 낡은 구건물과 정원수가 있는 토지를 ￦2,000,000에 현금 취득하였다. 20×7년부터 20×9년 말 사옥이 완공되었을 때까지 다음의 추가적인 현금지출과 수입이 발생하였다.

{\small\centering
\begin{tabular}{@{}l r@{}}
\hline
ㆍ토지 등기수수료 & ￦10,000 \\
ㆍ토지 취득세 & ￦20,000 \\
ㆍ낡은 구건물 철거비 & ￦40,000 \\
ㆍ낡은 구건물 철거로부터 발생한 & \\
\hspace{0.7em}골재 매각 순이익 & ￦10,000 \\
ㆍ신사옥에 대한 설계비 & ￦15,000 \\
ㆍ신사옥에 대한 공사대금 기지급액 & ￦2,000,000 \\
ㆍ토지위에 있는 정원수 처분으로 인한 순이익 & ￦270,000 \\
\hline
\end{tabular}\par}
20×9년 말 공사비 미지급액으로 ￦300,000이 남아있으며, 이 중에서 ￦50,000은 공사가 지체되었기 때문에 감액하기로 합의하였다. 이 경우에 토지와 건물의 취득원가는 각각 얼마인가? \par

{\small\centering
\begin{tabular}{@{}ccc@{}}
 & 토지 & 건물 \\
\cline{2-2}\cline{3-3}
① & ￦2,060,000 & ￦1,995,000 \\
② & ￦1,790,000 & ￦2,415,000 \\
③ & ￦1,800,000 & ￦1,995,000 \\
④ & ￦2,060,000 & ￦2,302,000 \\
⑤ & ￦1,790,000 & ￦2,265,000 \\
\end{tabular}\par}
\end{question}
\answer{5}
